% Independent audits: root and descent_groups; certificate code/C_finish_ramified_cancellation.py.
\begin{proposition}[Proved: cancellation of the rank-one witness in cyclotomic ramification]
\label{C:cyclotomic_torsion_cancellation}
Let $p$ be odd. Let $k$ be a commutative ring containing a unit
$\zeta$ with $\zeta^p=1$ and $1+\zeta+\cdots+\zeta^{p-1}=0$.
Let $I\subseteq k$ be an ideal with $I^2=0$ and
$(\zeta-1)I=0$. For $x\in I$, write $L_x$ for the free
$k[C_p]$-line on which the chosen generator acts by $1+x$, and put
$\delta(x)=[L_x]-[\mathbf1]\in R_k(C_p)$.
Then $x\mapsto\delta(x)$ is an additive homomorphism
$I\to R_k(C_p)$, and
\[
 \operatorname{Ind}_{C_p}^{\mathfrak S_p}\delta(x)=0
 \quad\text{for every }x\in I.
\]
In particular, if $A$ is a characteristic-zero DVR containing a
primitive $p$-th root of unity and $J$ is any $A$-module annihilated
by $\zeta-1$, the explicit rank-one induced obstruction of
Proposition~\ref{C:odd_torsion_allheights} vanishes over
$k=A\ltimes J$. This is not an assertion that $\vartheta_k$ is an
isomorphism.
\end{proposition}
\begin{proof}
Since $\zeta$ acts as one on $I$, its geometric-sum relation
implies $pI=0$. Put $z=[Z]\in R_k(C_p)$, where $Z$ is the line
with scalar action $\zeta$. Thus $z^p=1$ for the internal tensor
product.

For any $s\in I$, the matrix
\[
 T_s=\begin{pmatrix}1+s&1\\0&\zeta\end{pmatrix}
\]
has $T_s^p=1$. Its diagonal entries have $p$-th power one. Its
upper-right entry after taking the $p$-th power is
\[
 \sum_{i=0}^{p-1}(1+s)^i\zeta^{p-1-i}
 =\sum_{i=0}^{p-1}\zeta^{p-1-i}
   +s\sum_{i=0}^{p-1}i
 =s\frac{p(p-1)}2=0,
\]
where $p$ odd and $pI=0$ are used in the last equality.
For every $t\in I$, the invertible change of basis
\[
 P_{s,t}=\begin{pmatrix}1&0\\t-s&1\end{pmatrix}
\]
gives the exact identity
\[
 P_{s,t}^{-1}T_sP_{s,t}
   =\begin{pmatrix}1+t&1\\0&\zeta+s-t\end{pmatrix}.
\]
Both triangular presentations give admissible free-line exact
sequences. Because $\zeta(s-t)=s-t$, their equality in $R_k(C_p)$
is
\[
 \delta(s)=\delta(t)+z\delta(s-t).             \tag{C.1.29}
\]
Setting $s=0$ gives $\delta(t)=-z\delta(-t)$.
Applying this twice gives $z^2\delta(t)=\delta(t)$; since $p$
is odd and $z^p=1$, it follows that $z\delta(t)=\delta(t)$.
Equation~\textup{(C.1.29)} now proves
$\delta(s)=\delta(t)+\delta(s-t)$, as claimed.

Conjugation by the permutation $j\mapsto2j$ of $\mathbb F_p$
normalizes the standard $p$-cycle and raises its generator to the
second power. It identifies the inductions of $L_x$ and $L_{2x}$,
since $(1+x)^2=1+2x$. Therefore
\[
 \operatorname{Ind}\delta(x)=\operatorname{Ind}\delta(2x)
   =2\operatorname{Ind}\delta(x),
\]
and the induced class is zero. The last specialization satisfies
all the stated hypotheses, since the ideal $0\ltimes J$ is
square-zero and is killed by $\zeta-1$.
\end{proof}
